% chapters_parte2/04_datacenter_vxlan_evpn.tex

\chapter{\eh{office-building} Datacenter: VxLAN y EVPN}

Los datacenters modernos tienen un problema que las VLANs clásicas no pueden resolver:
millones de tenants (clientes), VMs que migran entre servidores físicos sin cambiar de IP,
y redes que deben estirarse entre edificios y regiones geográficas. VxLAN + EVPN es la
solución que usan AWS, Google Cloud, Azure y cualquier datacenter serio.

% ─────────────────────────────────────────────
\section{\eh{warning} Limitaciones de VLANs en Datacenter}

\begin{cuidado}
\begin{itemize}
    \item \textbf{4094 VLANs máximo.} Una nube pública tiene millones de clientes, cada
          uno necesita su red aislada. 4094 no alcanza ni para empezar.
    \item \textbf{STP no escala.} Con miles de switches, Spanning Tree converge lentamente
          y bloquea enlaces que podrían estar transmitiendo.
    \item \textbf{VMs migran físicamente.} Cuando vMotion mueve una VM de un servidor a
          otro, necesita mantener la misma IP (mismo L2). Con VLANs, eso requiere que
          ambos servidores estén en el mismo dominio L2, lo que limita la distancia física.
    \item \textbf{L2 sobre L3 es complejo.} Los datacenters rutean internamente (L3 escala),
          pero las aplicaciones esperan L2 (broadcast, ARP). VLANs no resuelven esto.
\end{itemize}
\end{cuidado}

% ─────────────────────────────────────────────
\section{\eh{package} VxLAN --- Virtual Extensible LAN}

\begin{concepto}
\textbf{VxLAN} crea un \emph{overlay} L2 sobre una red L3. Encapsula frames Ethernet
completos dentro de paquetes UDP, permitiendo que dos VMs en servidores físicamente
separados (incluso en datacenters distintos) crean que están en el mismo segmento Ethernet.

\begin{itemize}
    \item \textbf{VNI (VxLAN Network Identifier):} 24 bits $\to$ \textbf{16 millones}
          de segmentos virtuales. Resuelve el límite de 4094 VLANs.
    \item \textbf{VTEP (VxLAN Tunnel Endpoint):} el dispositivo que encapsula/desencapsula.
          Puede ser un switch físico (Leaf) o un servidor Linux con OVS.
    \item \textbf{Puerto UDP:} 4789 (estándar IANA).
\end{itemize}
\end{concepto}

\subsection{Formato del frame VxLAN}

\begin{lstlisting}
+---------------+----------+----------+-----------+------------------+
| Outer Ethernet| Outer IP | Outer UDP| VxLAN Hdr | Inner Ethernet   |
| (VTEP src/dst)|(VTEP IPs)| dst:4789 | (VNI 24b) | Frame original   |
| 14 bytes      | 20 bytes | 8 bytes  | 8 bytes   | variable         |
+---------------+----------+----------+-----------+------------------+
                  ^ Red física underlay              ^ Red virtual overlay
\end{lstlisting}

\begin{itemize}
    \item El frame Ethernet original (inner) viaja intacto dentro del túnel UDP.
    \item La red underlay solo ve IPs de VTEPs; no sabe nada de las VMs.
    \item \textbf{Overhead total:} 50 bytes. En enlaces de 10 GbE, es despreciable.
\end{itemize}

% ─────────────────────────────────────────────
\section{\eh{deciduous-tree} Topología Spine-Leaf}

\begin{concepto}
Las redes de datacenter tradicionales tienen 3 capas: Core $\to$ Distribution $\to$ Access.
Spanning Tree bloquea links redundantes y la latencia varía según el camino.

\textbf{Spine-Leaf} reemplaza esas 3 capas con solo 2, eliminando STP y maximizando
el uso de todos los enlaces mediante ECMP (Equal-Cost Multi-Path).
\end{concepto}

\subsection{Reglas del diseño Spine-Leaf}

\begin{itemize}
    \item Cada \textbf{Leaf} conecta a \textbf{todos} los Spines (ningún Leaf conecta
          directamente a otro Leaf).
    \item Cada \textbf{Spine} conecta a \textbf{todos} los Leaves (ningún Spine conecta
          a otro Spine).
    \item Entre cualquier par de servidores: \textbf{máximo 2 hops} (Leaf $\to$ Spine
          $\to$ Leaf). Latencia predecible y constante.
    \item Para escalar: agregar más Leaves (más puertos de servidor) o más Spines (más
          ancho de banda de east-west).
\end{itemize}

\subsection{Diagrama Spine-Leaf con VxLAN}

\begin{center}
\begin{tikzpicture}[
    spine/.style={draw=cyan!60!white, fill=darkcardviolet, text=textlight,
                  rectangle, rounded corners=4pt, minimum width=2.2cm,
                  minimum height=0.8cm, font=\small\bfseries, align=center},
    leaf/.style={draw=green!50!white, fill=darkcardgreen, text=textlight,
                 rectangle, rounded corners=4pt, minimum width=2.2cm,
                 minimum height=0.8cm, font=\small\bfseries, align=center},
    srv/.style={draw=gray!50!white, fill=darkcardgray, text=textlight,
                rectangle, minimum width=1.5cm, minimum height=0.6cm,
                font=\tiny, align=center},
    link/.style={thick, draw=gray!60!white},
    vtep/.style={text=yellow!70!white, font=\tiny\itshape},
]
    % Spines
    \node[spine] (s1) at (1.5,4) {Spine 1};
    \node[spine] (s2) at (4.5,4) {Spine 2};

    % Leaves
    \node[leaf] (l1) at (0,2) {Leaf 1\\VTEP};
    \node[leaf] (l2) at (3,2) {Leaf 2\\VTEP};
    \node[leaf] (l3) at (6,2) {Leaf 3\\VTEP};

    % Servers
    \node[srv] (srv1) at (-0.7,0.3) {VM-A\\10.0.1.10};
    \node[srv] (srv2) at (0.7,0.3)  {VM-B\\10.0.1.11};
    \node[srv] (srv3) at (2.3,0.3)  {VM-C\\10.0.2.10};
    \node[srv] (srv4) at (3.7,0.3)  {VM-D\\10.0.2.11};
    \node[srv] (srv5) at (5.3,0.3)  {VM-E\\10.0.3.10};
    \node[srv] (srv6) at (6.7,0.3)  {VM-F\\10.0.3.11};

    % Spine-Leaf links
    \draw[link] (l1) -- (s1); \draw[link] (l1) -- (s2);
    \draw[link] (l2) -- (s1); \draw[link] (l2) -- (s2);
    \draw[link] (l3) -- (s1); \draw[link] (l3) -- (s2);

    % Leaf-Server links
    \draw[link] (l1) -- (srv1); \draw[link] (l1) -- (srv2);
    \draw[link] (l2) -- (srv3); \draw[link] (l2) -- (srv4);
    \draw[link] (l3) -- (srv5); \draw[link] (l3) -- (srv6);

    % VxLAN tunnel annotation
    \draw[draw=yellow!60!white, dashed, thick, bend left=20]
        (l1) to node[vtep, above] {VxLAN tunnel VNI=100} (l3);
\end{tikzpicture}

{\small VM-A y VM-E están en el mismo VNI 100 (misma red virtual) aunque en Leaves distintos.
El tráfico viaja encapsulado en UDP: Leaf1 $\to$ Spine $\to$ Leaf3.}
\end{center}

% ─────────────────────────────────────────────
\section{\eh{satellite-antenna} EVPN --- Ethernet VPN}

\begin{concepto}
VxLAN sola tiene un problema: ¿cómo sabe un VTEP a qué otro VTEP enviar un frame si
no conoce la MAC destino? La respuesta naïve es \textbf{flood-and-learn}: inundar el
BUM traffic (Broadcast, Unknown-unicast, Multicast) a todos los VTEPs y aprender MACs
por el tráfico de respuesta. Eso no escala con miles de VTEPs.

\textbf{EVPN} usa \textbf{BGP} como protocolo de control para distribuir información
MAC/IP entre VTEPs. Cada VTEP anuncia las MACs/IPs de sus VMs locales vía BGP, y todos
los demás VTEPs aprenden esa información \emph{antes} de que llegue el primer paquete.
Resultado: no hay flooding innecesario.
\end{concepto}

\subsection{Route Types de EVPN}

\begin{center}
\begin{tabular}{lll}
\toprule
\textbf{Type} & \textbf{Nombre} & \textbf{Propósito} \\
\midrule
Type 1 & Ethernet A-D Route        & Multihoming, failover rápido \\
Type 2 & MAC/IP Advertisement      & Aprender MAC + IP de VMs remotas \\
Type 3 & Inclusive Multicast Route & Registro para recibir BUM traffic \\
Type 4 & Ethernet Segment Route    & Descubrimiento de multihoming \\
Type 5 & IP Prefix Route           & Routing L3 entre VNIs distintos \\
\bottomrule
\end{tabular}
\end{center}

\begin{ejemplo}
\textbf{Flujo con EVPN:}
\begin{enumerate}
    \item VM-A (en Leaf1) envía ARP Request buscando VM-E.
    \item Leaf1 tiene en su BGP EVPN table: VM-E (MAC=aa:bb:...) está en VTEP Leaf3
          (IP=10.1.1.3), VNI=100. Aprendido vía Type 2 Route.
    \item Leaf1 \textbf{no floodea}. Responde el ARP directamente (ARP suppression)
          con la MAC de VM-E desde su tabla local.
    \item VM-A envía el frame. Leaf1 lo encapsula en VxLAN hacia Leaf3.
    \item Leaf3 desencapsula y entrega a VM-E.
\end{enumerate}
Sin EVPN: el ARP hubiera sido floodeado a todos los Leaves. Con EVPN: cero flooding.
\end{ejemplo}

\subsection{Configuración básica (Cisco NX-OS)}

\begin{codigo}
\begin{lstlisting}
! Habilitar VxLAN y EVPN
feature nv overlay
feature bgp
nv overlay evpn

! NVE interface (VTEP)
interface nve1
  no shutdown
  source-interface loopback0
  member vni 10100
    ingress-replication protocol bgp

! BGP con address-family EVPN
router bgp 65001
  address-family l2vpn evpn
    advertise-pip

neighbor 10.1.1.2 remote-as 65001
  address-family l2vpn evpn
    send-community extended
\end{lstlisting}
\end{codigo}

% ─────────────────────────────────────────────
\section{\eh{up-arrow} Routing Asimétrico vs Simétrico en VxLAN}

\begin{center}
\begin{tabular}{lll}
\toprule
\textbf{Modelo} & \textbf{Dónde rutea} & \textbf{Trade-off} \\
\midrule
Asimétrico & VTEP de ingreso rutea, VTEP de egreso bridgea & Simple pero cada Leaf necesita todos los VNIs \\
Simétrico  & Ambos VTEPs rutean, L3 VNI dedicado           & Más escalable, cada Leaf solo necesita sus VNIs \\
\bottomrule
\end{tabular}
\end{center}

\begin{moderno}
\emoji{rocket} \textbf{VxLAN + EVPN en la nube pública:} AWS usa una variante propietaria
de estos principios internamente. Cada Availability Zone tiene su fabric de Leaf/Spine.
El tráfico entre AZs viaja encapsulado (propietario, no OpenFlow). Los Security Groups
de AWS se implementan sobre estos mismos mecanismos de overlay.
\end{moderno}
